% 第 48 章　两个量子系统相遇（完整样章）
\chapter{两个量子系统相遇}

到第~47~章为止，我们的量子账本一次只记一个系统：一个光子、一个自旋、一个量子比特。可真实世界里量子系统从不独居——它们碰撞、交换能量、被同一台仪器先后测量，然后各奔东西。本章的问题是：当两个量子系统相遇又分开之后，账本上会发生什么？答案里藏着量子力学最深、也最被滥用的一个概念：纠缠。它能让两个粒子隔着上千公里保持一种任何经典安排都伪造不出来的默契；它也让无数科幻故事兴高采烈地宣布“超光速通信实现了”。本章的任务是把这两件事都落到实处：前者是真事，我们会算出它究竟“超”在哪里；后者是误读，我们会精确地指出它错在哪一行。

\step{反常识开场}

先从一件你今天就能做的事说起。2016~年起，IBM~把真实的量子计算机挂到了互联网上，任何人注册一个免费账号就能远程使用。最简单的入门程序只有几步：取两个量子比特，施加两个标准操作（一个制造叠加，一个让两者相互作用，名字暂且不管），然后分别测量这两个比特，重复一千次。把结果抄成两列：第一列是第一个比特的一千个结果，第二列是第二个比特的。

现在做两件事。第一件：捂住第二列，只看第一列。你会看到世界上最彻底的随机序列——0~和~1~大约各占一半，前后没有任何规律，统计学家用尽办法也挑不出毛病。第二件：把手拿开，把两列并排对齐。你会愣住：两列几乎逐位相同。第一列是~0~的地方，第二列几乎总是~0；是~1~的地方几乎总是~1。一千对结果，失配的只有零星几个，而那还可以归罪于机器噪声。

请认真掂量这件事有多怪：每一列单独看，是教科书级的完美噪声；两列合起来看，是教科书级的严格同步。一枚完美混乱的硬币和另一枚完美混乱的硬币，掷出了一模一样的结果，连着一千次。它们是怎么“约好”的？

更麻烦的还在后面。这样的粒子对并不只存在于那台机器里：用一束激光照射一种叫~BBO~的特殊晶体，晶体中会成对地冒出光子，每对光子之间就有这种“各自随机、合起来同步”的关系。物理学家早已把这样的光子对分开送往远处——先是实验室两端，后来是几十公里的光纤，再后来是卫星与地面之间——相关依然分毫不差。而且我们还可以让两个粒子先飞远，再临时决定用哪种方式测量：它们连“商量对策”的时间都没有。

先把两个最容易糊弄的“解释”挡在门外。一个说：这不过是两台机器都被造得很准而已。可你家里的两台电子设备，哪怕同厂同批，各自的随机数也绝不会逐位对上；制造得再像的两枚硬币，掷一万次也对不上几次。另一个说：它们之间有某种我们没发现的无线电联系。可把探测器裹进屏蔽罩、搬到地下室、送上卫星，相关一丝不减；而且后面会看到，任何“联系”的说法都要先过一道数学关卡，而这道关卡恰恰是它过不去的。本章结束时，你会知道这种默契到底是什么、它为什么不是普通的“两个结果相关”、以及它为什么依然寄不出一封超光速的信。

\step{先猜一猜}

在继续之前，把你的猜测写下来，越具体越好：

第一个猜测是关于机制的。（a）每个比特内部有个随机发生器，但两个比特在制备时被装进了同一份程序——就像工厂把两张号码相同的彩票封进两个信封，相关性来自“出生证明”。（b）它们出生时彼此独立，但测量其中一个的瞬间，它向另一个发出了某种信号：“我被测了，结果是~0，你也选~0。”这个信号或许极快，甚至无限快。（c）两个比特从头到尾就不是“两个东西”，而是某一个整体的两半；问“每个比特各自是什么状态”，这个问题本身就不合法。

第二个猜测是关于应用的。把这样一对粒子分给相隔一千公里的两个人。甲方通过“测”与“不测”向乙方发摩尔斯电码：这小时测量，表示“点”；不测，表示“划”。乙方只需要盯着自己那列随机数。你猜，乙方能不能分辨出甲方“测了没有”？

第三个猜测是关于限度的。任何“出生时就写好答案”的经典安排，是否存在一个理论上限，而量子粒子对可以突破它？如果存在，你敢下注突破多少——百分之一？百分之十？还是根本没有上限，想多强就多强？

把三份答案夹在书页里。读完本章回来对账，看看你的直觉输在哪一步。

\step{动手实验}

\begin{experiment}[贝尔游戏：亲手摸到“出生证明”的天花板]
\textbf{材料}：一个朋友，两枚硬币（或两部能出随机数的手机），纸和笔，屋子里的两个角落。

\textbf{规则}：每一轮，你们各自抛硬币，决定自己被问到哪个问题：正面为问题~1，反面为问题~2。然后各自在纸上写下一个答案：0~或~1。评分规则只有一条：如果两人都被问了问题~2，你们赢的条件是答案\textbf{不同}；其余三种提问组合（至少一人被问问题~1），赢的条件是答案\textbf{相同}。

\textbf{玩法}：开局之前，你们可以商量任何策略——“永远答~0”“被问问题~1~答~0、问题~2~答~1”，甚至现场掷骰子决定答案，都可以。但每轮抛硬币之后，不许再看对方、不许再说话。玩四十轮，记下赢了多少轮。

\textbf{任务}：想办法赢过~75\%，也就是四十轮里赢三十轮以上。
\end{experiment}

你会失败的，而且败得很讲理。最好的策略是“两个人永远答相同的数”：四种提问组合里，它能赢下三种，只在“双双被问问题~2”时输掉——恰好~75\%。换一个策略也一样：任何预先写好的答案表，在四种组合里最多只能满足三种。掷骰子随机答？那不过是在几张答案表之间抽签，赢率还是封顶~75\%。这是算术，不是技巧问题。

这个游戏的正式名字叫~CHSH~游戏（四位物理学家姓氏的首字母）。它还有下半截：如果允许两人各带一个纠缠粒子，每轮用事先约定的方向测量自己的粒子来产生答案，量子力学给出的获胜率是~\(\cos^2(\pi/8)\approx85\%\)。注意，粒子在游戏开始后同样没有任何交流。你在客厅里用尽办法也翻不过去的~75\%，被一对光子安静地越过——实验室里早已反复实现。这就是开场悬念的量化版本。

值得停下来体会这个游戏的设计意图。它把“有没有预先写好的答案表”这个看似纯哲学的问题，变成了一道可以动手批改的算术题：四种提问组合是考卷，答案表是考生的底牌，75\%~是任何底牌都逃不掉的分数上限。爱因斯坦一方与玻尔一方争论了三十年“粒子被测量之前到底有没有确定属性”，谁也说服不了谁；贝尔的贡献，是发现这场争论里藏着一个能被实验判决的预言差。哲学问题一旦被翻译成数字，就有了裁判。

还有一个顺手的观察。让你和朋友各自用手机生成一长串随机数，对齐比较：相同位数的比例大约是~50\%，毫无相关。这说明“随机”本身遍地都是，一点也不稀奇；稀奇的是“各自纯随机，合起来严丝合缝”。

\step{旧模型遇到麻烦}

先给猜测~(a)~一个体面的名字：局域隐变量模型，大白话叫“出生证明”理论。1935~年，爱因斯坦、波多尔斯基和罗森的著名论文（EPR）把它推到了最锋利的形式。他们的推理是：量子力学声称粒子对分开后各自没有确定状态，可测量结果又完美相关——那唯一合理的解释是，每个粒子从出生起就揣着一份完整的答案表，对任何可能的测量方向都写好了答案；相关，是因为两张表一开始就互相配合。量子力学之所以看不见这份表，是因为它“不完备”，像一本只印了摘要的书。这个模型有两个不可分割的零件：其一，答案预先存在；其二，分开之后互不打扰，甲方怎么测不影响乙方。

“出生证明”理论确实能解释完美相关：工厂装两份相同的程序就行。它的死穴在于：当可选的测量方式不止一种时，它必须对每一轮游戏交出一份预先写好的完整答案表——而你的家庭实验刚刚证明，任何完整答案表在贝尔游戏里最多赢~75\%。1964~年，贝尔把这个常识观察变成了严格定理：凡满足“预先答案加互不打扰”的模型，在多组测量方向下的相关强度有一个数学上限，这就是贝尔不等式。而量子力学对纠缠态的预言，越过这个上限大约十五个百分点。这不是哲学口味的分歧，是一道算术题，两边不可能同时都对。

那就只剩猜测~(b)：测量瞬间发信号。这个模型有两处致命伤。其一，实验不给面子：人们让两个探测器相距很远，把测量压缩在极短的时间窗里完成，快到任何信号得以远超光速的若干倍飞行才来得及——相关依然成立，而且数值依然恰好是量子力学的预言，不多不少。其二，理论不允许：量子力学可以严格证明，无论甲方对自己的粒子做什么操作，乙方单看自己那列数据的统计分布都纹丝不动。这叫无信号定理。换句话说，即便大自然内部真在用某种超光速的机制“对表”，它也同时保证这种机制永远兑现不成一封能寄出去的信。这个“即便……也……”的结构很微妙，我们到模型的边界一节再回来掂量。

现在算一笔带真实数据的账，免得上面的“很远”“极短”听起来像修辞。2015~年，荷兰代尔夫特理工的小组把两颗纠缠的电子自旋放在相距~1.3~公里的两个实验室里。光走完这段路需要
\[
\frac{1.3\times10^{3}\ \text{m}}{3\times10^{8}\ \text{m/s}}\approx4.3\times10^{-6}\ \text{s},
\]
而两边的测量从开始到落槌只用了约一微秒——任何不超光速的“商量”都来不及发生。那一轮实验测得的贝尔参数约为~2.42，明显高于经典上限~2（低于量子理论上限~\(2\sqrt2\approx2.83\)~的部分可由损耗和噪声解释）。同年，美国和奥地利的小组用光子得到同样结论，堵上了最后几个技术性漏洞。此后，纠缠分发被推到上千公里：2017~年，中国的“墨子号”卫星把纠缠光子对发到相距约~1200~公里的两个地面站，相关依然违反贝尔不等式。2022~年的诺贝尔物理学奖授予了克劳泽、阿斯佩和塞林格，表彰的正是这条从思想实验一路走到无漏洞实验的战线。“出生证明”理论不是被投票淘汰的，是被测距仪和秒表淘汰的。

\begin{figure}[htbp]
\centering
\begin{tikzpicture}[>=Stealth,scale=1.0]
  \node[draw,circle,fill=orange!40,minimum size=9mm] (S) at (0,0) {源};
  \node[draw,fill=blue!15,minimum width=24mm,minimum height=11mm,align=center] (A) at (-4.6,0) {甲实验室\\可转测量方向};
  \node[draw,fill=blue!15,minimum width=24mm,minimum height=11mm,align=center] (B) at (4.6,0) {乙实验室\\可转测量方向};
  \draw[->,thick] (S) -- (A) node[midway,above]{光子 1};
  \draw[->,thick] (S) -- (B) node[midway,above]{光子 2};
  \draw[<->,dashed] (-4.6,-1.3) -- (4.6,-1.3)
    node[midway,below,align=center]{1.3 km：光要飞 4.3 \(\mu\)s，测量只花约 1 \(\mu\)s\\“商量”来不及发生};
\end{tikzpicture}
\caption{贝尔实验的要点不在设备精巧，而在时间安排：两边测量落槌之快，使任何不超光速的串供都来不及完成。相关依旧，且超过经典上限。}
\end{figure}

\step{发明新概念}

旧账本的毛病出在哪一页？出在一条我们从没审过的默认假设上：“联合系统的状态等于两个子系统状态的合并清单”。量子力学的新规矩是：联合系统的可能性不是相加，而是相乘。一个比特有~2~个可能状态；两个比特有~\(2\times2=4\)~个；十个比特有~1024~个。这套“各选各的、组合计数”的构造叫张量积——名字唬人，内容就是小学乘法。第~46~章说过，量子态要写成所有可能性的叠加，于是两个比特的完整账本有四栏：00、01、10、11，每栏配一个振幅。

关键一步来了。在这张四栏账本里，大多数账目可以拆成“甲的账本乘乙的账本”：比如“甲必为~0、乙一半一半”，各写各的，互不牵连，这叫可分态。但有一类账目怎么拆也拆不开，最典型的是
\[
|\Phi^{+}\rangle=\frac{1}{\sqrt2}\big(|00\rangle+|11\rangle\big).
\]
请直接朗读它：“同为~0”与“同为~1”各占一半，01~和~10~两栏是空的。注意这个账目的奇特性格：对整体，它写得清清楚楚、没有一丝含糊——整个系统处于一个完全确定的状态；可如果你只问“甲这个比特是什么状态”，答案是“一半~0~一半~1，纯随机”。整体完全确定，每个部分却毫无确定性——这种组合在经典世界里根本不存在：你知道一双手套的全部，当然知道每一只手套。这种拆不开的联合态，就叫纠缠态。

张量积这条“小学乘法”还有一个吓人的后果，值得在这里立一块路标。栏目数随比特数翻倍地涨：十个比特一千栏，一百个比特约~\(10^{30}\)~栏，三百个比特约~\(10^{90}\)~栏——超过可观测宇宙的原子总数。换句话说，一个中等规模的量子系统的完整账本，是用全宇宙的原子当纸也抄不下的。这不是修辞，而是“量子计算机为什么可能强大”的真正起点：大自然处理这本巨账毫不费力，经典计算机却连记账都记不起。至于这笔财富为什么很难取出来花，挑战题里再细算。

纠缠不是什么，必须一条条钉死。它不是“两个粒子之间有根看不见的线”；不是“它们一直在互相通信”；也不是“两个结果碰巧一致”。同卵双胞胎的考试分数接近，是因为共享了基因和书桌——那是经典相关，答案事前各有定数；纠缠粒子在被问到之前，连“各自的定数”都不存在，存在的只有整体的那一笔账。再用比喻划一次界：它像一双分装两盒的手套，因为打开一只立刻知道另一只；它不像手套，因为手套的颜色在装箱时就已写定，而纠缠粒子的相关强度超过任何“装箱时写定”之安排的上限——你的贝尔游戏已经替这条界限站过岗。

新概念的第一个定理级推论是不可克隆定理：一个未知的量子态不能被完美复制。复印机可以复印文件，因为文件上的每个比特都能先读再写；但“先读”一个量子比特等于对它做一次测量，测量只回答你选定的那一个问题，还顺手把状态改写掉——你永远拿不到完整的原件，自然印不出第二份。证明只有三行，放在数学桥里；这里只记住它的地位：不可克隆不是眼下的工程瓶颈，而是线性叠加结构的直接后果，定律级别，与能量守恒同级。

\step{一句话模型}

\begin{oneline}
两个量子系统一旦相互作用过，就可以处于一种只能整体书写、拆不成两份独立档案的联合态——纠缠态：对整体，账目完全确定；对每个部分单独看，是纯随机；把两边的测量记录事后对齐，相关强度超过任何“出生时写好答案”的经典安排；但这种相关永远兑现不成一个比特的超光速信息。
\end{oneline}

这句话像什么？像两份逐字对应、却各自用一次性密码本加密的电报：单独看任何一份都是噪声，两份合在一起才显出内容。它不像什么？不像心灵感应——“感应”意味着一边的念头流到了另一边，而这里没有任何东西在两地之间流动；也不像被同一根弹簧拴住的两颗小球，因为小球各自每时每刻都有确定的位置，而纠缠粒子的部分连“确定的状态”都没有。

\step{数学翻译器}

先在主线层做两笔账。第一笔，单边账目。纠缠对~\(|\Phi^{+}\rangle\) 里，甲测得~0~的概率等于~\(|00\rangle\) 项与~\(|01\rangle\) 项的概率之和：\(\tfrac12+0=\tfrac12\)；乙同理。两列数据各自完美随机——开场的两列噪声有了出处。第二笔，对齐账目。“甲得~0~且乙得~0”的概率是~\(\tfrac12\)，“甲得~1~且乙得~1”也是~\(\tfrac12\)，交叉组合为零：两列必然逐位相同。做特殊情形检查：如果叠加中~\(|11\rangle\) 的系数缩成~0，账目退化成“两人都是~0”的普通事实，随机性和纠缠一起消失——对，系数为零正是“可分”与“纠缠”的分水岭。

实验里更常用的是自旋版纠缠对。用第~47~章的语言，两个自旋可以组成单态
\[
|\Psi^{-}\rangle=\frac{1}{\sqrt2}\big(|{\uparrow\downarrow}\rangle-|{\downarrow\uparrow}\rangle\big).
\]
甲乙沿同一方向测量，结果必然相反；若两人选的测量方向夹角为~\(\theta\)，大量重复后两边结果的相关程度是
\[
E(\theta)=-\cos\theta .
\]
按惯例回答四个问题。它描述什么：每次测量两边各得~\(+1\)~或~\(-1\)，\(E(\theta)\)~是大量重复后两个结果乘积的平均值。为什么这样写：\(\cos\theta\)~是两个方向之间几何关系的唯一出场方式，负号来自单态“总和为零”的构造。何时成立：理想制备、理想测量、退相干可忽略。何时失效：真实实验有损耗和噪声，实测值的绝对值会缩小——代尔夫特的~2.42~与~\(2\sqrt2\)~之间的差距正是这么来的。特殊角盘问：\(\theta=0^\circ\)~时~\(E=-1\)，完全反相关，对；\(\theta=90^\circ\)~时~\(E=0\)，互相垂直的提问彼此独立；\(\theta=180^\circ\)~时~\(E=+1\)，反着测等于翻个面，完全相关。三个锚点都符合直觉，而中间角度才是经典模型跟不上的地方。

\begin{figure}[htbp]
\centering
\begin{tikzpicture}[>=Stealth,scale=0.9]
  % 1 比特
  \node at (0,2.6) {1 个比特：2 栏};
  \draw[fill=blue!15] (-0.8,1.6) rectangle (0.8,2.2); \node at (0,1.9) {0};
  \draw[fill=blue!15] (-0.8,0.8) rectangle (0.8,1.4); \node at (0,1.1) {1};
  % 2 比特
  \node at (3.4,2.6) {2 个比特：4 栏};
  \foreach \y/\t in {1.6/00, 0.8/01, 0.0/10, -0.8/11}{
    \draw[fill=blue!25] (2.6,\y) rectangle (4.2,\y+0.6); \node at (3.4,\y+0.3) {\t};
  }
  % 3 比特
  \node at (6.8,2.6) {3 个比特：8 栏};
  \foreach \y/\t in {2.0/000, 1.3/001, 0.6/010, -0.1/011, -0.8/100, -1.5/101, -2.2/110, -2.9/111}{
    \draw[fill=blue!35] (6.0,\y) rectangle (7.6,\y+0.55); \node at (6.8,\y+0.27) {\footnotesize \t};
  }
  \node[below] at (3.4,-1.5) {可能性相乘：\(2\to4\to8\)};
  \node[below] at (6.8,-3.6) {每栏配一个振幅，才是完整账本};
\end{tikzpicture}
\caption{张量积的直觉：联合系统的栏目数是各系统栏目数的乘积，不是和。纠缠态就是一张“拆不成两份小账本”的大账本。}
\end{figure}

\begin{mathbridge}
\textbf{“拆不开”的代数判决。}试试把~\(|\Phi^{+}\rangle\) 硬拆成两份小账本。两个比特各自的一般状态是~\(a|0\rangle+b|1\rangle\) 与~\(c|0\rangle+d|1\rangle\)，相乘展开得
\[
(ac)|00\rangle+(ad)|01\rangle+(bc)|10\rangle+(bd)|11\rangle .
\]
要等于~\(\tfrac{1}{\sqrt2}(|00\rangle+|11\rangle)\)，必须~\(ac=bd=\tfrac{1}{\sqrt2}\)（都不为零），同时~\(ad=bc=0\)。但~\(ad=0\)~要求~\(a\)~或~\(d\)~为零，无论哪个都会让~\(ac\)~或~\(bd\)~变成零——矛盾。四个方程无解，所以拆不开。可分与纠缠，在代数上就是“四个系数能不能写成两两乘积”这一件事。

\textbf{不可克隆定理的三行证明。}假设存在一台通用复制机，能把任何输入态~\(|\psi\rangle\) 变成两份~\(|\psi\rangle|\psi\rangle\)。它对~\(|0\rangle\) 应输出~\(|00\rangle\)，对~\(|1\rangle\) 应输出~\(|11\rangle\)。量子演化是线性的，所以输入叠加态~\(\tfrac{1}{\sqrt2}(|0\rangle+|1\rangle)\) 时，输出只能是两个输出的叠加~\(\tfrac{1}{\sqrt2}(|00\rangle+|11\rangle)\)——一个纠缠态，而不是我们想要的~\(\tfrac12(|0\rangle+|1\rangle)(|0\rangle+|1\rangle)\)。要逐条正确，线性性不答应；服从线性性，叠加态就复制错。结论：任何机器都复制不了完全未知的量子态。

\textbf{贝尔游戏的~85\%~从哪来。}CHSH~游戏的常见写法是定义一个组合量~\(S\)，对应正文里“按四种提问组合的输赢计分”。任何经典策略满足~\(S\leq2\)，换算成赢率恰是~75\%；量子力学让两人按问题选择夹角合适的测量方向，可得~\(S=2\sqrt2\)，换算赢率~\(\cos^2(\pi/8)\approx85\%\)。还有一个常被忽略的优雅事实：\(2\sqrt2\)~也是量子理论自身的上限——连大自然自己玩这个游戏，也赢不到~100\%。“相关能多强”这件事，既有地板也有天花板。
\end{mathbridge}

\begin{challenge}
\textbf{CHSH~不等式一行推导。}设甲方对两种提问的“预先答案”为~\(a,a'\in\{+1,-1\}\)，乙方为~\(b,b'\)。考察组合
\[
S=ab+ab'+a'b-a'b'=a(b+b')+a'(b-b').
\]
无论~\(b,b'\)~取何值，\((b+b')\)~与~\((b-b')\)~必有一个是~\(\pm2\)、另一个是~0，所以每一轮~\(S=\pm2\)，任何“预先答案”模型的平均值都满足~\(|S|\leq2\)。量子力学对单态选取两组夹角~\(45^\circ\)~的测量方向，可以算出~\(|S|=2\sqrt2>2\)。实验的判决书由此而来：只要测得~\(S>2\)，就宣判“局域”与“预先答案”这对联合假设死刑。注意判决书的主语是联合假设——想单独抢救哪一个，各有各的代价，而这正是诠释分歧的入口。

\textbf{部分是纯随机的正式说法。}“只看甲这一半”的正式工具叫约化密度矩阵。对~\(|\Phi^{+}\rangle\) 算出甲的约化态，结果是~\(\rho_A=\tfrac12\big(|0\rangle\langle0|+|1\rangle\langle1|\big)\)——最大混合态，对任何测量方向都给出五五开。无信号定理在此变成一行矩阵运算：乙方无论测不测、沿什么方向测，甲方的~\(\rho_A\)~都分毫不动。纠缠的定量度量（纠缠熵）也正是对这个约化态计算的：整体越纯、部分越混，纠缠越深。单态是“最纠缠”的两个比特——它的每一半都是最大程度的无知。
\end{challenge}

\step{模型的边界}

照例把本章内容分成三堆。实验事实：贝尔不等式在几十年间上百个实验里被稳定违反，2015~年后的版本已无技术漏洞；单边数据永远是纯随机，无信号从未被打破；纠缠可以制备、分发、维持到上千公里，并随环境干扰而衰减。数学模型：张量积结构加纠缠态加玻恩规则，预言精度极高。它的适用条件：非相对论性情境、粒子数固定、系统与环境隔离得足够好。失效的方向也有路标：退相干会把纠缠“漏”给环境，宏观物体几乎立刻失去纠缠的可观测后果；粒子数可以产生和湮灭时，张量积语言要升级成量子场论（第~49~章）。物理解释：对“测量落槌瞬间，另一边发生了什么”，所有诠释的共同预言只有一句——“单边统计毫无变化”；至于整体账目的更新是不是某种真实的瞬时事件，哥本哈根、多世界、玻姆、QBism~各讲各的故事，现有实验一视同仁。任何人告诉你“实验证明了超光速影响”，他都是把诠释当成了事实。

再列一张典型误用清单。“量子通信等于超光速通信”——错，量子密钥分发和隐形传态都必须借助普通经典信道核对，速度不超光速。“纠缠粒子上刻着另一半的全部信息，可以隔空取物”——错，单边账目是纯随机，什么也刻不上去。“两个人的意识通过纠缠相连”——账本里从头到尾没有意识这个栏目。“量子计算机就是~\(2^{300}\)~台经典电脑并行干活”——过度简化，错在哪里留给挑战题。

最后补一条通向前后章节的桥。本章只讨论了“两个干净的系统”之间的纠缠；可只要把乙换成“环境”——空气分子、热辐射、仪器里亿万个自由度——你立刻得到第~45~章退相干的新说法：所谓退相干，就是系统与环境缠在了一起，那本只能整体书写的账，被摊薄到再也收不回来的无数栏目里。纠缠本身没有消失，只是扩散到了你无法对齐、也无法核对的角落。这解释了为什么纠缠在实验室里娇贵无比：工程师与噪声作战，本质上是在阻止系统同环境“共享账本”。反过来说，本章的每一个结论也都是第~45~章测量问题的延续——测量仪器与被测系统之间建立的，正是一段精心设计、又能被读出来的纠缠。

\step{回头解释开场现象}

回到云端那两列数据，现在每一行都有了账目。制备阶段，两个比特被做成~\(|\Phi^{+}\rangle\)：整体账目确定，各自账目纯随机——这解释了“每列都像完美噪声”。测量后对齐：00~与~11~各占一半，交叉项为零——这解释了“逐位相同”。把装置搬远、把测量方向留到最后一刻临时改选，相关依旧，而且稳稳越过~75\%~的游戏上限——于是“出厂程序”被判死刑；“隔空发信”既拿不出收据，也兑不了现。剩下的正是本章的模型：那不是两台约好了的硬币，而是一枚同时写在两个地址上的硬币。

开场猜的摩尔斯电码方案，现在可以亲手拆穿。甲方“测或不测”，乙方手里的序列都是同样深度的噪声，统计上不可区分——无信号定理封死了这条路。甲方哪怕更换测量方向，改变的也只是事后两列对齐时的相关花样，乙方单列依旧一片盲然。想读出任何信息，必须等甲方把记录用普通方式寄过来对齐——电话、光纤、信鸽，统统不超光速。纠缠是“事后才能核实的默契”，不是“即时送达的电报”。

顺带一提，这套机制还给出了一份意想不到的礼物：可被公证的随机数。你电脑里的“随机数”其实大多是按固定算法算出来的伪随机，懂行的人原则上能预测；而纠缠对给出的随机性带着贝尔不等式的背书——只要测得的相关超过了经典上限，就证明这些结果不是任何“预先写好的程序”的产物，随机性是自然定律认证的。这类“设备无关”的随机数发生器已经从论文走进了产品，算是贝尔那场哲学官司最实在的赔偿。

那这种兑不了现的默契有什么用？两个已经在运行的答案。其一，量子密钥分发：双方分发纠缠对、各自随机选方向测量，再公开一小部分记录，核对贝尔型相关是否还在。任何窃听者中途插手测量，都会把纠缠拆坏，相关强度当场跌破阈值——窃听行为在统计上无处遁形。中国的“京沪干线”已经把这条原理跑在两千公里的光纤网络上。其二，量子隐形传态：要把一个未知量子态从甲处搬到乙处，办法是让甲乙共享一对纠缠粒子，甲对自己的粒子和待传的态做一次联合测量，把两个比特的结果用经典信道发给乙，乙据此对自己的粒子做一个对应操作——原件在甲处被测量毁掉，副本在乙处生成，不可克隆定理全程安然无恙，经典信道卡着光速的闸门。2017~年，“墨子号”把隐形传态从地面送上了千里之外的卫星轨道。请注意两个应用的共同点：纠缠负责“经典物理给不出的关联”，光速负责“信息依然得一步一步走”。

\step{脑洞实验与挑战题}

\begin{thought}[把账本写到星系两端]
设想把一对纠缠光子分别发往相距一光年的两座空间站，约定明年的同一天正午同时测量。结果依旧严格相关。有人问：这一年的飞行里，它们“记得”彼此吗？本章的账本回答：不存在两个“它们”各自记着什么；记得一切的是那一份同时写在两个地址上的联合态。再问：甲站站长赶在正午前临时改主意、换一个测量方向，能给乙站“留个言”吗？不能——乙站正午拿到的记录，与甲站测不测、怎么测，在统计上毫无关系；相关要等双方把记录用光速电波凑到一起时才显现。还有一个更让人不舒服的问题：如果整体账目在两次测量落槌时才“更新”，这次更新横跨一光年，它算“发生”在哪一刻？第~39~章说过，不同运动状态的观察者连“哪次测量在先”都可以意见不一。好在无信号定理保证：无论你把这笔账记在哪个时刻，任何可观测的后果都不依赖这个选择。量子力学与相对论之所以能和平共处，靠的就是这种“内部账本可以非局域、外部后果必须局域”的分工。至于这本内部账本的性质，正是第~49~章把量子力学和相对论缝在一起时要重新翻开的一页。
\end{thought}

\begin{puzzles}
\item 严格证明你在贝尔游戏里赢不过~75\%。把“被问问题~1~答什么、问题~2~答什么”写成答案表（两人共四个数），论证：四种提问组合里至少有一种必输。\emph{提示}：两人答案始终相同，能赢三种组合、输掉“双问题~2”。想让“双问题~2”也赢，两人答案就必须不同，这要求恰好一人按问题切换答案；逐个检查“恰有一人被问问题~2”的两种组合，会发现总有一个填不平。随机化只是在几张答案表之间抽签，救不了场。
\item 有人说：“不可克隆只是今天的仪器不行，未来总能造出量子复印机。”用本章内容反驳，并说明：假如量子复印机真的存在，会连锁引发什么后果？\emph{提示}：复印机加纠缠对再加两套测量基，等于一台超光速电报机——乙方把甲方发来的单个粒子复制一万份，沿两个方向各测一半，就能读出甲方当初选了哪套基。所以不可克隆与无信号是互相担保的两块砖：抽掉一块，整座量子大厦的因果结构就塌了。
\item 你的云端数据里混进了~5\%~的失配。朋友说：“看来纠缠也没多神，照样出错。”请区分两种“出错”：噪声造成的失配，与“甲方换测量方向导致乙方序列改变”，并说明为什么前者天天发生、后者从未被观测到。\emph{提示}：失配只在两列对齐时才现身；乙方单列的全部统计（0~与~1~的比例、任何内部规律）都与甲方的行为无关。“纠缠被环境拆坏”与“信息被传了过来”，要用完全不同的对照实验来检验。
\item 三百个量子比特的账本有~\(2^{300}\approx10^{90}\)~栏，超过可观测宇宙的原子总数（约~\(10^{80}\)）。于是有人说：“量子计算机等于~\(10^{90}\)~台电脑同时开工。”指出这句话对在哪、错在哪。\emph{提示}：对的部分——经典计算机确实记不下这本账，这正是量子计算诱人的原因；错的部分——一次测量只吐出某一栏的答案，其余栏目的“工作”你根本读不到。量子算法的全部艺术，是让正确答案的栏目在测量前因干涉而涨大、错误答案相互抵消。这本巨账不是劳动力，是乐谱。
\end{puzzles}
